\documentclass[11pt,a4paper]{article}

%=========================================================
% PACKAGES
%=========================================================
\usepackage[utf8]{inputenc}
\usepackage[T1]{fontenc}
\usepackage[french]{babel}
\usepackage[a4paper,margin=2.2cm]{geometry}
\usepackage{lmodern}
\usepackage{amsmath,amssymb,amsthm,mathtools}
\usepackage{fancyhdr}
\usepackage{lastpage}
\usepackage{enumitem}
\usepackage{array,tabularx}
\usepackage{tikz}
\usepackage[most]{tcolorbox}
\usepackage{hyperref}
\usepackage{xcolor}

%=========================================================
% MISE EN PAGE
%=========================================================
\setlength{\headheight}{14pt}
\pagestyle{fancy}
\fancyhf{}
\fancyhead[L]{Mathématiques CPGE}
\fancyhead[C]{Réduction d'endomorphismes}
\fancyhead[R]{Page \thepage/\pageref{LastPage}}
\fancyfoot[C]{Cours complet -- Algèbre linéaire}

\hypersetup{
    colorlinks=true,
    linkcolor=blue!50!black,
    urlcolor=blue!50!black
}

\setlist[itemize]{label=\(\triangleright\)}
\setlist[enumerate]{label=\textbf{\arabic*.}}

%=========================================================
% COULEURS ET BOÎTES
%=========================================================
\definecolor{bleufonce}{RGB}{25,65,130}
\definecolor{bleuclair}{RGB}{235,243,255}
\definecolor{vertfonce}{RGB}{30,120,70}
\definecolor{vertclair}{RGB}{235,250,240}
\definecolor{rougefonce}{RGB}{170,40,40}
\definecolor{rougeclair}{RGB}{255,238,238}
\definecolor{orangefonce}{RGB}{190,110,20}
\definecolor{orangeclair}{RGB}{255,246,230}
\definecolor{violetfonce}{RGB}{95,55,145}
\definecolor{violetclair}{RGB}{246,241,255}

\tcbset{
    enhanced,
    sharp corners,
    boxrule=0.7pt,
    before skip=8pt,
    after skip=8pt
}

\newtcolorbox{definitionbox}[1][]{
    colback=bleuclair,
    colframe=bleufonce,
    title=\textbf{Définition #1}
}

\newtcolorbox{theorembox}[1][]{
    colback=vertclair,
    colframe=vertfonce,
    title=\textbf{Théorème / Propriété #1}
}

\newtcolorbox{methodbox}[1][]{
    colback=orangeclair,
    colframe=orangefonce,
    title=\textbf{Méthode #1}
}

\newtcolorbox{retenirbox}[1][]{
    colback=violetclair,
    colframe=violetfonce,
    title=\textbf{À retenir #1}
}

\newtcolorbox{erreurbox}[1][]{
    colback=rougeclair,
    colframe=rougefonce,
    title=\textbf{Erreur classique #1}
}

\newtcolorbox{exercicebox}[1][]{
    colback=white,
    colframe=black!55,
    title=\textbf{Exercice #1}
}

\newtcolorbox{correctionbox}[1][]{
    colback=gray!5,
    colframe=black!65,
    title=\textbf{Correction #1}
}

%=========================================================
% COMMANDES
%=========================================================
\newcommand{\K}{\mathbb K}
\newcommand{\R}{\mathbb R}
\newcommand{\C}{\mathbb C}
\newcommand{\N}{\mathbb N}
\newcommand{\M}{\mathcal M}
\newcommand{\Lcal}{\mathcal L}
\newcommand{\Vect}{\operatorname{Vect}}
\newcommand{\Ker}{\operatorname{Ker}}
\newcommand{\Imf}{\operatorname{Im}}
\newcommand{\rg}{\operatorname{rg}}
\newcommand{\tr}{\operatorname{tr}}
\newcommand{\Id}{\operatorname{Id}}
\newcommand{\Sp}{\operatorname{Sp}}
\newcommand{\diag}{\operatorname{diag}}

\newtheorem{prop}{Proposition}[section]
\newtheorem{lemme}[prop]{Lemme}
\newtheorem{corollaire}[prop]{Corollaire}

%=========================================================
% DOCUMENT
%=========================================================
\begin{document}

\begin{center}
    {\Huge \textbf{Réduction d'endomorphismes}}\\[2mm]
    {\Large Valeurs propres, polynôme caractéristique, diagonalisation et trigonalisabilité}\\[2mm]
    {\large Cours complet pour classes préparatoires scientifiques}\\[2mm]
    \rule{0.88\linewidth}{0.5pt}
\end{center}

\vspace{0.4cm}

\tableofcontents

\newpage

%=========================================================
\section*{Introduction générale}
\addcontentsline{toc}{section}{Introduction générale}

Soit \(E\) un espace vectoriel de dimension finie sur un corps \(\K\), généralement \(\R\) ou \(\C\), et soit
\[
u\in\Lcal(E).
\]
L'objectif de la réduction est de comprendre l'action de \(u\) en choisissant une base aussi adaptée que possible.

Dans une base quelconque, la matrice de \(u\) peut être compliquée.  
Mais dans une bonne base, elle peut devenir :
\begin{itemize}
    \item diagonale ;
    \item triangulaire ;
    \item ou au moins plus simple à étudier.
\end{itemize}

\begin{retenirbox}
Réduire un endomorphisme, c'est choisir une base qui révèle sa structure.
\[
\text{Endomorphisme compliqué}
\quad\longrightarrow\quad
\text{matrice simple dans une bonne base.}
\]
\end{retenirbox}

La diagonalisation est la situation idéale. Si
\[
A=PDP^{-1},
\]
avec \(D\) diagonale, alors :
\[
A^n=PD^nP^{-1},
\]
et les puissances de \(D\) se calculent immédiatement :
\[
D^n=\diag(\lambda_1^n,\dots,\lambda_p^n).
\]

La réduction intervient dans de nombreux domaines :
\begin{itemize}
    \item calcul des puissances de matrices ;
    \item résolution de suites récurrentes vectorielles ;
    \item systèmes différentiels linéaires ;
    \item chaînes de Markov ;
    \item géométrie linéaire ;
    \item étude des projecteurs, symétries et endomorphismes particuliers.
\end{itemize}

%=========================================================
\section{Pré-requis issus des applications linéaires}

\subsection{Applications linéaires et endomorphismes}

\begin{definitionbox}[Application linéaire]
Soient \(E\) et \(F\) deux espaces vectoriels sur \(\K\).  
Une application
\[
u:E\to F
\]
est linéaire si :
\[
\forall x,y\in E,\ \forall \lambda,\mu\in\K,
\qquad
u(\lambda x+\mu y)=\lambda u(x)+\mu u(y).
\]
\end{definitionbox}

\begin{definitionbox}[Endomorphisme]
Si \(E=F\), une application linéaire
\[
u:E\to E
\]
est appelée un \textbf{endomorphisme} de \(E\).  
On note l'ensemble des endomorphismes de \(E\) :
\[
\Lcal(E).
\]
\end{definitionbox}

\begin{definitionbox}[Automorphisme]
Un endomorphisme bijectif est appelé un \textbf{automorphisme}.
\end{definitionbox}

\begin{retenirbox}
La réduction concerne les endomorphismes, donc les applications linéaires d'un espace vectoriel dans lui-même.
\end{retenirbox}

\subsection{Noyau, image et rang}

\begin{definitionbox}[Noyau et image]
Soit \(u:E\to F\) linéaire.  
Le noyau de \(u\) est :
\[
\Ker(u)=\{x\in E:\ u(x)=0_F\}.
\]
L'image de \(u\) est :
\[
\Imf(u)=\{u(x):x\in E\}.
\]
\end{definitionbox}

\begin{theorembox}
Le noyau \(\Ker(u)\) est un sous-espace vectoriel de \(E\), et l'image \(\Imf(u)\) est un sous-espace vectoriel de \(F\).
\end{theorembox}

\begin{proof}
On montre d'abord que le noyau est stable par combinaison linéaire.  
Soient \(x,y\in\Ker(u)\) et \(\lambda,\mu\in\K\). Alors :
\[
u(\lambda x+\mu y)=\lambda u(x)+\mu u(y)=\lambda 0_F+\mu 0_F=0_F.
\]
Donc \(\lambda x+\mu y\in\Ker(u)\).

Pour l'image, si \(z_1,z_2\in\Imf(u)\), alors il existe \(x_1,x_2\in E\) tels que :
\[
z_1=u(x_1),\qquad z_2=u(x_2).
\]
Alors :
\[
\lambda z_1+\mu z_2
=
\lambda u(x_1)+\mu u(x_2)
=
u(\lambda x_1+\mu x_2)\in\Imf(u).
\]
\end{proof}

\begin{definitionbox}[Rang]
Si \(E\) est de dimension finie, le rang de \(u\) est :
\[
\rg(u)=\dim(\Imf(u)).
\]
\end{definitionbox}

\begin{theorembox}[Théorème du rang]
Si \(u:E\to F\) est linéaire et si \(E\) est de dimension finie, alors :
\[
\dim(E)=\dim(\Ker(u))+\rg(u).
\]
\end{theorembox}

\begin{retenirbox}
Le théorème du rang est indispensable pour calculer les dimensions des sous-espaces propres :
\[
E_\lambda(u)=\Ker(u-\lambda\Id_E).
\]
\end{retenirbox}

\subsection{Matrices et changement de base}

\begin{definitionbox}[Matrice d'un endomorphisme]
Soit \(E\) de dimension \(n\), et soit
\[
\mathcal B=(e_1,\dots,e_n)
\]
une base de \(E\).  
La matrice de \(u\in\Lcal(E)\) dans la base \(\mathcal B\), notée
\[
\operatorname{Mat}_{\mathcal B}(u),
\]
est la matrice dont la \(j\)-ième colonne est le vecteur des coordonnées de \(u(e_j)\) dans la base \(\mathcal B\).
\end{definitionbox}

\begin{theorembox}[Changement de base]
Soient \(\mathcal B\) et \(\mathcal B'\) deux bases de \(E\).  
Soit \(P\) la matrice de passage de \(\mathcal B'\) vers \(\mathcal B\), c'est-à-dire la matrice dont les colonnes sont les coordonnées des vecteurs de \(\mathcal B'\) exprimés dans \(\mathcal B\).

Si :
\[
A=\operatorname{Mat}_{\mathcal B}(u),
\qquad
A'=\operatorname{Mat}_{\mathcal B'}(u),
\]
alors :
\[
A'=P^{-1}AP.
\]
\end{theorembox}

\begin{proof}
Soit \(x\in E\).  
On note \(X\) le vecteur colonne des coordonnées de \(x\) dans \(\mathcal B\), et \(X'\) celui dans \(\mathcal B'\).  
Par définition de \(P\), on a :
\[
X=PX'.
\]
Dans la base \(\mathcal B\), les coordonnées de \(u(x)\) sont :
\[
AX=APX'.
\]
Pour revenir dans la base \(\mathcal B'\), on multiplie par \(P^{-1}\).  
Les coordonnées de \(u(x)\) dans \(\mathcal B'\) sont donc :
\[
P^{-1}APX'.
\]
Ainsi :
\[
A'=P^{-1}AP.
\]
\end{proof}

\begin{definitionbox}[Matrices semblables]
Deux matrices \(A,B\in\M_n(\K)\) sont dites semblables s'il existe \(P\in GL_n(\K)\) tel que :
\[
B=P^{-1}AP.
\]
\end{definitionbox}

\begin{retenirbox}
Deux matrices semblables représentent le même endomorphisme dans deux bases différentes.
\end{retenirbox}

\begin{methodbox}[Passer du point de vue géométrique au point de vue matriciel]
\begin{itemize}
    \item Un endomorphisme \(u\) correspond à une matrice \(A\) dans une base.
    \item Changer de base revient à remplacer \(A\) par une matrice semblable \(P^{-1}AP\).
    \item Diagonaliser \(u\), c'est trouver une base où sa matrice est diagonale.
    \item Diagonaliser \(A\), c'est trouver \(P\) inversible telle que \(P^{-1}AP\) soit diagonale.
\end{itemize}
\end{methodbox}

%=========================================================
\section{Valeurs propres, vecteurs propres et sous-espaces propres}

\subsection{Définitions}

\begin{definitionbox}[Valeur propre et vecteur propre]
Soit \(u\in\Lcal(E)\).  
Un scalaire \(\lambda\in\K\) est une \textbf{valeur propre} de \(u\) s'il existe un vecteur non nul \(x\in E\) tel que :
\[
u(x)=\lambda x.
\]
Un tel vecteur \(x\) est appelé \textbf{vecteur propre} associé à la valeur propre \(\lambda\).
\end{definitionbox}

\begin{erreurbox}
Le vecteur nul n'est jamais un vecteur propre.  
En effet, il vérifie toujours \(u(0)=\lambda 0\), quelle que soit la valeur de \(\lambda\), ce qui ne donne aucune information.
\end{erreurbox}

\begin{definitionbox}[Spectre]
L'ensemble des valeurs propres de \(u\) est appelé le \textbf{spectre} de \(u\), et se note :
\[
\Sp(u).
\]
Pour une matrice \(A\), on note de même :
\[
\Sp(A).
\]
\end{definitionbox}

\begin{definitionbox}[Sous-espace propre]
Soit \(\lambda\in\K\).  
Le sous-espace propre associé à \(\lambda\) est :
\[
E_\lambda(u)=\Ker(u-\lambda\Id_E).
\]
Autrement dit :
\[
E_\lambda(u)=\{x\in E:\ u(x)=\lambda x\}.
\]
\end{definitionbox}

\begin{theorembox}
Pour tout \(\lambda\in\K\), \(E_\lambda(u)\) est un sous-espace vectoriel de \(E\).
\end{theorembox}

\begin{proof}
On a :
\[
E_\lambda(u)=\Ker(u-\lambda\Id_E).
\]
Comme \(u-\lambda\Id_E\) est une application linéaire, son noyau est un sous-espace vectoriel de \(E\).
\end{proof}

\begin{theorembox}
Un scalaire \(\lambda\) est une valeur propre de \(u\) si et seulement si :
\[
E_\lambda(u)\neq\{0\}.
\]
\end{theorembox}

\begin{proof}
Par définition, \(\lambda\) est une valeur propre s'il existe \(x\neq0\) tel que :
\[
u(x)=\lambda x.
\]
Cela équivaut à :
\[
(u-\lambda\Id_E)(x)=0.
\]
Donc \(x\in\Ker(u-\lambda\Id_E)=E_\lambda(u)\), avec \(x\neq0\).  
Ainsi :
\[
E_\lambda(u)\neq\{0\}.
\]
La réciproque est immédiate.
\end{proof}

\begin{retenirbox}
Les vecteurs propres associés à \(\lambda\), auxquels on ajoute le vecteur nul, forment exactement le sous-espace propre :
\[
E_\lambda(u).
\]
\end{retenirbox}

\subsection{Traduction matricielle}

Soit \(A\in\M_n(\K)\).  
Un scalaire \(\lambda\) est valeur propre de \(A\) s'il existe une colonne non nulle \(X\in\M_{n,1}(\K)\) telle que :
\[
AX=\lambda X.
\]
Cela équivaut à :
\[
(A-\lambda I_n)X=0.
\]

Donc :
\[
E_\lambda(A)=\Ker(A-\lambda I_n).
\]

\begin{methodbox}[Calculer un sous-espace propre]
Pour calculer \(E_\lambda(A)\) :
\begin{enumerate}
    \item écrire la matrice \(A-\lambda I_n\) ;
    \item résoudre le système homogène
    \[
    (A-\lambda I_n)X=0 ;
    \]
    \item donner une base du noyau obtenu ;
    \item la dimension de ce noyau est la multiplicité géométrique de \(\lambda\).
\end{enumerate}
\end{methodbox}

%=========================================================
\section{Premières propriétés des valeurs propres}

\subsection{Homothéties}

\begin{theorembox}
Les valeurs propres de l'homothétie
\[
u=\lambda\Id_E
\]
sont réduites à \(\lambda\), si \(E\neq\{0\}\).
\end{theorembox}

\begin{proof}
Soit \(\mu\) une valeur propre de \(u\).  
Il existe \(x\neq0\) tel que :
\[
u(x)=\mu x.
\]
Or :
\[
u(x)=\lambda x.
\]
Donc :
\[
\lambda x=\mu x,
\]
d'où :
\[
(\lambda-\mu)x=0.
\]
Comme \(x\neq0\), on obtient :
\[
\lambda-\mu=0,
\]
donc \(\mu=\lambda\).
\end{proof}

\subsection{Matrices triangulaires}

\begin{theorembox}
Les valeurs propres d'une matrice triangulaire sont ses coefficients diagonaux.
\end{theorembox}

\begin{proof}
Soit \(A=(a_{ij})\) une matrice triangulaire supérieure.  
Alors :
\[
XI_n-A
\]
est encore triangulaire supérieure, de coefficients diagonaux :
\[
X-a_{11},\dots,X-a_{nn}.
\]
Donc :
\[
\det(XI_n-A)=\prod_{k=1}^n (X-a_{kk}).
\]
Les racines sont donc exactement les coefficients diagonaux de \(A\).
\end{proof}

\begin{erreurbox}
Une matrice triangulaire n'est pas forcément diagonalisable.  
Elle permet de lire les valeurs propres, mais pas de conclure automatiquement à la diagonalisation.
\end{erreurbox}

%=========================================================
\section{Polynôme caractéristique}

\subsection{Définition}

\begin{definitionbox}[Polynôme caractéristique d'une matrice]
Soit \(A\in\M_n(\K)\).  
Le polynôme caractéristique de \(A\) est :
\[
\chi_A(X)=\det(XI_n-A).
\]
\end{definitionbox}

\begin{definitionbox}[Polynôme caractéristique d'un endomorphisme]
Si \(u\in\Lcal(E)\), avec \(E\) de dimension finie, on définit :
\[
\chi_u(X)=\chi_A(X),
\]
où \(A\) est la matrice de \(u\) dans une base quelconque de \(E\).
\end{definitionbox}

\begin{theorembox}
Le polynôme caractéristique d'un endomorphisme ne dépend pas de la base choisie.
\end{theorembox}

\begin{proof}
Soient \(A\) et \(B\) les matrices de \(u\) dans deux bases différentes.  
Alors \(A\) et \(B\) sont semblables : il existe \(P\in GL_n(\K)\) tel que :
\[
B=P^{-1}AP.
\]
Alors :
\[
XI_n-B
=
XI_n-P^{-1}AP
=
P^{-1}(XI_n-A)P.
\]
Donc :
\[
\det(XI_n-B)
=
\det(P^{-1})\det(XI_n-A)\det(P).
\]
Comme :
\[
\det(P^{-1})\det(P)=1,
\]
on obtient :
\[
\det(XI_n-B)=\det(XI_n-A).
\]
Donc :
\[
\chi_B=\chi_A.
\]
\end{proof}

\subsection{Valeurs propres et polynôme caractéristique}

\begin{theorembox}
Un scalaire \(\lambda\in\K\) est une valeur propre de \(A\in\M_n(\K)\) si et seulement si :
\[
\chi_A(\lambda)=0.
\]
\end{theorembox}

\begin{proof}
On a :
\[
\lambda\in\Sp(A)
\]
si et seulement s'il existe \(X\neq0\) tel que :
\[
AX=\lambda X.
\]
Cela équivaut à :
\[
(A-\lambda I_n)X=0.
\]
Donc \(A-\lambda I_n\) n'est pas inversible, ce qui équivaut à :
\[
\det(A-\lambda I_n)=0.
\]
Or :
\[
\det(\lambda I_n-A)=(-1)^n\det(A-\lambda I_n).
\]
Donc :
\[
\det(\lambda I_n-A)=0.
\]
Autrement dit :
\[
\chi_A(\lambda)=0.
\]
\end{proof}

\begin{retenirbox}
Les valeurs propres sont exactement les racines du polynôme caractéristique.
\[
\lambda\in\Sp(A)
\Longleftrightarrow
\chi_A(\lambda)=0.
\]
\end{retenirbox}

\subsection{Trace, déterminant et valeurs propres}

\begin{theorembox}
Si le polynôme caractéristique de \(A\in\M_n(\K)\) est scindé :
\[
\chi_A(X)=\prod_{k=1}^n (X-\lambda_k),
\]
où les valeurs propres sont répétées avec multiplicité, alors :
\[
\tr(A)=\lambda_1+\cdots+\lambda_n,
\]
et :
\[
\det(A)=\lambda_1\cdots\lambda_n.
\]
\end{theorembox}

\begin{proof}
On écrit :
\[
\chi_A(X)=\det(XI_n-A).
\]
Le coefficient de \(X^{n-1}\) dans ce polynôme est \(-\tr(A)\).  
D'autre part :
\[
\prod_{k=1}^n(X-\lambda_k)
=
X^n-(\lambda_1+\cdots+\lambda_n)X^{n-1}+\cdots.
\]
En identifiant les coefficients de \(X^{n-1}\), on obtient :
\[
\tr(A)=\lambda_1+\cdots+\lambda_n.
\]

En évaluant en \(X=0\), on a :
\[
\chi_A(0)=\det(-A)=(-1)^n\det(A).
\]
Mais :
\[
\chi_A(0)=\prod_{k=1}^n(-\lambda_k)=(-1)^n\lambda_1\cdots\lambda_n.
\]
Donc :
\[
\det(A)=\lambda_1\cdots\lambda_n.
\]
\end{proof}

%=========================================================
\section{Familles de vecteurs propres}

\subsection{Lemme fondamental}

\begin{theorembox}[Liberté des vecteurs propres]
Des vecteurs propres associés à des valeurs propres deux à deux distinctes sont linéairement indépendants.
\end{theorembox}

\begin{proof}
Soient \(\lambda_1,\dots,\lambda_p\) des valeurs propres deux à deux distinctes, et soient
\[
x_1,\dots,x_p
\]
des vecteurs propres non nuls associés respectivement à ces valeurs propres.

On montre que la famille \((x_1,\dots,x_p)\) est libre par récurrence sur \(p\).

Pour \(p=1\), la famille \((x_1)\) est libre car \(x_1\neq0\).

Supposons le résultat vrai au rang \(p-1\).  
Considérons une relation linéaire :
\[
\alpha_1x_1+\cdots+\alpha_px_p=0.
\]
En appliquant \(u\), on obtient :
\[
\alpha_1u(x_1)+\cdots+\alpha_pu(x_p)=0,
\]
donc :
\[
\alpha_1\lambda_1x_1+\cdots+\alpha_p\lambda_px_p=0.
\]
Multiplions la première relation par \(\lambda_p\) :
\[
\alpha_1\lambda_px_1+\cdots+\alpha_p\lambda_px_p=0.
\]
En soustrayant, on obtient :
\[
\alpha_1(\lambda_1-\lambda_p)x_1+\cdots+
\alpha_{p-1}(\lambda_{p-1}-\lambda_p)x_{p-1}=0.
\]
Par hypothèse de récurrence, la famille \((x_1,\dots,x_{p-1})\) est libre.  
Donc, pour tout \(i\leq p-1\),
\[
\alpha_i(\lambda_i-\lambda_p)=0.
\]
Comme \(\lambda_i\neq\lambda_p\), on a :
\[
\alpha_i=0.
\]
La relation initiale devient :
\[
\alpha_px_p=0.
\]
Comme \(x_p\neq0\), on obtient :
\[
\alpha_p=0.
\]
Tous les coefficients sont nuls : la famille est libre.
\end{proof}

\begin{corollaire}
Les sous-espaces propres associés à des valeurs propres deux à deux distinctes sont en somme directe.
\end{corollaire}

\begin{proof}
Soient \(\lambda_1,\dots,\lambda_r\) des valeurs propres distinctes.  
Pour montrer que
\[
E_{\lambda_1}\oplus\cdots\oplus E_{\lambda_r}
\]
est une somme directe, on considère une relation :
\[
x_1+\cdots+x_r=0,
\]
où \(x_i\in E_{\lambda_i}\).  
Si certains \(x_i\) étaient non nuls, ils formeraient une famille de vecteurs propres associés à des valeurs propres distinctes, donc libre.  
La seule relation possible est donc :
\[
x_1=\cdots=x_r=0.
\]
\end{proof}

\begin{corollaire}
Si \(E\) est de dimension \(n\) et si \(u\) admet \(n\) valeurs propres deux à deux distinctes, alors \(u\) est diagonalisable.
\end{corollaire}

\begin{proof}
Choisissons un vecteur propre non nul associé à chacune des \(n\) valeurs propres.  
D'après le théorème précédent, on obtient une famille libre de \(n\) vecteurs dans un espace de dimension \(n\).  
C'est donc une base de \(E\), formée de vecteurs propres.  
Ainsi \(u\) est diagonalisable.
\end{proof}

%=========================================================
\section{Multiplicités algébrique et géométrique}

\subsection{Définitions}

\begin{definitionbox}[Multiplicité algébrique]
Soit \(\lambda\in\Sp(u)\).  
La multiplicité algébrique de \(\lambda\), notée \(m_\lambda\), est sa multiplicité comme racine du polynôme caractéristique :
\[
\chi_u(X).
\]
\end{definitionbox}

\begin{definitionbox}[Multiplicité géométrique]
La multiplicité géométrique de \(\lambda\) est :
\[
\dim(E_\lambda(u)).
\]
\end{definitionbox}

\begin{theorembox}[Inégalité fondamentale]
Pour toute valeur propre \(\lambda\),
\[
1\leq \dim(E_\lambda(u))\leq m_\lambda.
\]
\end{theorembox}

\begin{proof}
Comme \(\lambda\) est une valeur propre, on a :
\[
E_\lambda(u)\neq\{0\}.
\]
Donc :
\[
\dim(E_\lambda(u))\geq1.
\]

Posons :
\[
r=\dim(E_\lambda(u)).
\]
Choisissons une base :
\[
(e_1,\dots,e_r)
\]
de \(E_\lambda(u)\), que l'on complète en une base de \(E\) :
\[
(e_1,\dots,e_r,e_{r+1},\dots,e_n).
\]
Pour \(1\leq i\leq r\), on a :
\[
u(e_i)=\lambda e_i.
\]
Dans cette base, la matrice de \(u\) est donc de la forme :
\[
A=
\begin{pmatrix}
\lambda I_r & B\\
0 & C
\end{pmatrix}.
\]
Alors :
\[
XI_n-A=
\begin{pmatrix}
(X-\lambda)I_r & -B\\
0 & XI_{n-r}-C
\end{pmatrix}.
\]
Le déterminant d'une matrice triangulaire par blocs vaut le produit des déterminants diagonaux. Donc :
\[
\chi_u(X)
=
\det(XI_n-A)
=
(X-\lambda)^r\det(XI_{n-r}-C).
\]
Ainsi \((X-\lambda)^r\) divise \(\chi_u(X)\).  
Donc :
\[
r\leq m_\lambda.
\]
\end{proof}

\begin{erreurbox}
Si l'on trouve :
\[
\dim(E_\lambda)>m_\lambda,
\]
alors le calcul est faux.  
Cela ne peut jamais arriver.
\end{erreurbox}

%=========================================================
\section{Diagonalisation}

\subsection{Matrices diagonales}

\begin{definitionbox}[Matrice diagonale]
Une matrice diagonale est une matrice de la forme :
\[
D=
\begin{pmatrix}
\lambda_1&0&\cdots&0\\
0&\lambda_2&\cdots&0\\
\vdots&\vdots&\ddots&\vdots\\
0&0&\cdots&\lambda_n
\end{pmatrix}.
\]
On note :
\[
D=\diag(\lambda_1,\dots,\lambda_n).
\]
\end{definitionbox}

\begin{retenirbox}
Si \(D=\diag(\lambda_1,\dots,\lambda_n)\), alors :
\[
D^k=\diag(\lambda_1^k,\dots,\lambda_n^k).
\]
\end{retenirbox}

\subsection{Définition de la diagonalisation}

\begin{definitionbox}[Endomorphisme diagonalisable]
Un endomorphisme \(u\in\Lcal(E)\) est diagonalisable s'il existe une base de \(E\) formée de vecteurs propres de \(u\).
\end{definitionbox}

\begin{definitionbox}[Matrice diagonalisable]
Une matrice \(A\in\M_n(\K)\) est diagonalisable s'il existe \(P\in GL_n(\K)\) et une matrice diagonale \(D\) telles que :
\[
A=PDP^{-1}.
\]
De manière équivalente :
\[
P^{-1}AP=D.
\]
\end{definitionbox}

\begin{retenirbox}
Dans la relation :
\[
A=PDP^{-1},
\]
les colonnes de \(P\) sont des vecteurs propres de \(A\), et les coefficients diagonaux de \(D\) sont les valeurs propres correspondantes, dans le même ordre.
\end{retenirbox}

\subsection{Critère fondamental}

\begin{theorembox}[Critère géométrique de diagonalisation]
Un endomorphisme \(u\in\Lcal(E)\) est diagonalisable si et seulement si :
\[
E=\bigoplus_{\lambda\in\Sp(u)} E_\lambda(u).
\]
\end{theorembox}

\begin{proof}
Supposons \(u\) diagonalisable.  
Il existe une base de \(E\) formée de vecteurs propres de \(u\).  
En regroupant les vecteurs associés à une même valeur propre, on obtient une base de chaque sous-espace propre, et la réunion de ces bases forme une base de \(E\).  
Donc :
\[
E=\bigoplus_{\lambda\in\Sp(u)}E_\lambda(u).
\]

Réciproquement, supposons :
\[
E=\bigoplus_{\lambda\in\Sp(u)}E_\lambda(u).
\]
Pour chaque valeur propre \(\lambda\), choisissons une base de \(E_\lambda(u)\).  
Comme la somme est directe et égale à \(E\), la réunion de ces bases est une base de \(E\).  
Tous les vecteurs de cette base sont des vecteurs propres.  
Donc \(u\) est diagonalisable.
\end{proof}

\begin{corollaire}
Un endomorphisme \(u\) est diagonalisable si et seulement si :
\[
\sum_{\lambda\in\Sp(u)}\dim(E_\lambda(u))=\dim(E).
\]
\end{corollaire}

\subsection{Critère par les multiplicités}

\begin{theorembox}[Critère pratique de diagonalisation]
Un endomorphisme \(u\in\Lcal(E)\), avec \(E\) de dimension \(n\), est diagonalisable sur \(\K\) si et seulement si :
\begin{enumerate}
    \item son polynôme caractéristique \(\chi_u\) est scindé sur \(\K\) ;
    \item pour toute valeur propre \(\lambda\),
    \[
    \dim(E_\lambda(u))=m_\lambda.
    \]
\end{enumerate}
\end{theorembox}

\begin{proof}
Supposons \(u\) diagonalisable.  
Dans une base de vecteurs propres, la matrice de \(u\) est diagonale :
\[
D=\diag(\lambda_1,\dots,\lambda_n).
\]
Donc :
\[
\chi_u(X)=\prod_{k=1}^n(X-\lambda_k),
\]
ce qui montre que \(\chi_u\) est scindé.

De plus, pour chaque valeur propre \(\lambda\), la multiplicité algébrique \(m_\lambda\) est le nombre d'apparitions de \(\lambda\) sur la diagonale, et la dimension de \(E_\lambda(u)\) est également ce nombre.  
Donc :
\[
\dim(E_\lambda(u))=m_\lambda.
\]

Réciproquement, supposons que \(\chi_u\) soit scindé et que :
\[
\forall\lambda\in\Sp(u),\quad \dim(E_\lambda(u))=m_\lambda.
\]
Comme \(\chi_u\) est scindé de degré \(n\), la somme des multiplicités algébriques des valeurs propres vaut \(n\).  
Donc :
\[
\sum_{\lambda\in\Sp(u)}\dim(E_\lambda(u))
=
\sum_{\lambda\in\Sp(u)}m_\lambda
=
n.
\]
Les sous-espaces propres étant en somme directe, leur somme a dimension \(n\), donc elle est égale à \(E\).  
Ainsi \(u\) est diagonalisable.
\end{proof}

\begin{methodbox}[Méthode complète de diagonalisation]
Pour diagonaliser une matrice \(A\in\M_n(\K)\) :
\begin{enumerate}
    \item Calculer le polynôme caractéristique :
    \[
    \chi_A(X)=\det(XI_n-A).
    \]
    \item Factoriser \(\chi_A\) sur le corps de base \(\K\).
    \item Déterminer les valeurs propres.
    \item Pour chaque valeur propre \(\lambda\), calculer :
    \[
    E_\lambda(A)=\Ker(A-\lambda I_n).
    \]
    \item Comparer :
    \[
    \dim(E_\lambda)
    \quad\text{et}\quad
    m_\lambda.
    \]
    \item Si la somme des dimensions des sous-espaces propres vaut \(n\), choisir une base de chaque sous-espace propre.
    \item Former la matrice \(P\) avec ces vecteurs propres en colonnes.
    \item Former \(D\) en mettant les valeurs propres correspondantes sur la diagonale.
    \item Écrire :
    \[
    A=PDP^{-1}.
    \]
\end{enumerate}
\end{methodbox}

%=========================================================
\section{Exemples détaillés de diagonalisation}

\subsection{Exemple 1 : deux valeurs propres distinctes}

Soit :
\[
A=
\begin{pmatrix}
1&1\\
0&2
\end{pmatrix}.
\]

On calcule :
\[
\chi_A(X)
=
\det
\begin{pmatrix}
X-1&-1\\
0&X-2
\end{pmatrix}
=
(X-1)(X-2).
\]
Les valeurs propres sont :
\[
1\quad\text{et}\quad 2.
\]
Elles sont distinctes. Comme \(A\in\M_2(\R)\), \(A\) est diagonalisable.

Calculons les sous-espaces propres.

Pour \(\lambda=1\) :
\[
A-I_2=
\begin{pmatrix}
0&1\\
0&1
\end{pmatrix}.
\]
Le système donne :
\[
y=0.
\]
Donc :
\[
E_1=\Vect
\begin{pmatrix}
1\\0
\end{pmatrix}.
\]

Pour \(\lambda=2\) :
\[
A-2I_2=
\begin{pmatrix}
-1&1\\
0&0
\end{pmatrix}.
\]
Le système donne :
\[
-y+x=0,
\]
donc \(y=x\). Ainsi :
\[
E_2=\Vect
\begin{pmatrix}
1\\1
\end{pmatrix}.
\]

On prend :
\[
P=
\begin{pmatrix}
1&1\\
0&1
\end{pmatrix},
\qquad
D=
\begin{pmatrix}
1&0\\
0&2
\end{pmatrix}.
\]
Alors :
\[
A=PDP^{-1}.
\]

\subsection{Exemple 2 : polynôme scindé mais matrice non diagonalisable}

Soit :
\[
A=
\begin{pmatrix}
1&1\\
0&1
\end{pmatrix}.
\]

On a :
\[
\chi_A(X)=(X-1)^2.
\]
L'unique valeur propre est \(1\), de multiplicité algébrique \(2\).

On calcule :
\[
A-I_2=
\begin{pmatrix}
0&1\\
0&0
\end{pmatrix}.
\]
Le système :
\[
(A-I_2)X=0
\]
donne :
\[
y=0.
\]
Donc :
\[
E_1=\Vect
\begin{pmatrix}
1\\0
\end{pmatrix}.
\]
Ainsi :
\[
\dim(E_1)=1<2=m_1.
\]
Donc \(A\) n'est pas diagonalisable.

\begin{erreurbox}
Un polynôme caractéristique scindé ne suffit pas à garantir la diagonalisation.  
Il faut aussi que, pour chaque valeur propre :
\[
\dim(E_\lambda)=m_\lambda.
\]
\end{erreurbox}

%=========================================================
\section{Polynômes d'endomorphismes}

\subsection{Définition}

\begin{definitionbox}[Polynôme d'un endomorphisme]
Soit \(u\in\Lcal(E)\) et soit :
\[
P(X)=a_0+a_1X+\cdots+a_mX^m\in\K[X].
\]
On définit :
\[
P(u)=a_0\Id_E+a_1u+\cdots+a_mu^m.
\]
\end{definitionbox}

\begin{definitionbox}[Polynôme d'une matrice]
Pour \(A\in\M_n(\K)\), on définit :
\[
P(A)=a_0I_n+a_1A+\cdots+a_mA^m.
\]
\end{definitionbox}

\subsection{Propriétés algébriques}

\begin{theorembox}
Pour tous \(P,Q\in\K[X]\), on a :
\[
(P+Q)(u)=P(u)+Q(u),
\]
\[
(PQ)(u)=P(u)\circ Q(u)=Q(u)\circ P(u).
\]
\end{theorembox}

\begin{proof}
La première égalité vient directement de l'addition terme à terme.

Pour le produit, il suffit de vérifier sur les monômes :
\[
X^iX^j=X^{i+j}.
\]
Donc :
\[
u^i\circ u^j=u^{i+j}.
\]
Par bilinéarité du produit de polynômes, on obtient :
\[
(PQ)(u)=P(u)\circ Q(u).
\]
Comme les puissances de \(u\) commutent entre elles, on a aussi :
\[
P(u)\circ Q(u)=Q(u)\circ P(u).
\]
\end{proof}

\subsection{Valeurs propres et polynômes}

\begin{theorembox}
Si \(x\) est vecteur propre de \(u\) associé à \(\lambda\), alors :
\[
P(u)(x)=P(\lambda)x.
\]
\end{theorembox}

\begin{proof}
Comme :
\[
u(x)=\lambda x,
\]
on montre par récurrence que :
\[
u^k(x)=\lambda^k x.
\]
Donc :
\[
P(u)(x)
=
a_0x+a_1u(x)+\cdots+a_mu^m(x)
=
(a_0+a_1\lambda+\cdots+a_m\lambda^m)x
=
P(\lambda)x.
\]
\end{proof}

\begin{definitionbox}[Polynôme annulateur]
Un polynôme \(P\in\K[X]\) est un polynôme annulateur de \(u\) si :
\[
P(u)=0.
\]
\end{definitionbox}

\begin{theorembox}
Si \(P(u)=0\), alors toute valeur propre \(\lambda\) de \(u\) vérifie :
\[
P(\lambda)=0.
\]
\end{theorembox}

\begin{proof}
Soit \(x\neq0\) un vecteur propre associé à \(\lambda\).  
On a :
\[
0=P(u)(x)=P(\lambda)x.
\]
Comme \(x\neq0\), cela impose :
\[
P(\lambda)=0.
\]
\end{proof}

\begin{erreurbox}
Les racines d'un polynôme annulateur contiennent les valeurs propres possibles.  
Mais toutes les racines du polynôme annulateur ne sont pas forcément des valeurs propres.
\end{erreurbox}

%=========================================================
\section{Diagonalisation par polynômes annulateurs}

\subsection{Théorème de décomposition des noyaux}

\begin{theorembox}[Décomposition des noyaux]
Soient \(P_1,\dots,P_r\in\K[X]\) deux à deux premiers entre eux.  
Soit :
\[
P=P_1P_2\cdots P_r.
\]
Si \(u\in\Lcal(E)\), alors :
\[
\Ker(P(u))
=
\Ker(P_1(u))\oplus\cdots\oplus\Ker(P_r(u)).
\]
\end{theorembox}

\begin{proof}
On donne la preuve pour deux polynômes, le cas général s'obtenant par récurrence.

Soient \(P,Q\in\K[X]\) premiers entre eux.  
Par le théorème de Bézout, il existe \(A,B\in\K[X]\) tels que :
\[
AP+BQ=1.
\]
En évaluant en \(u\), on obtient :
\[
A(u)P(u)+B(u)Q(u)=\Id_E.
\]

Soit \(x\in\Ker((PQ)(u))\).  
Alors :
\[
P(u)Q(u)(x)=0.
\]
On écrit :
\[
x=A(u)P(u)(x)+B(u)Q(u)(x).
\]
Posons :
\[
x_Q=A(u)P(u)(x),
\qquad
x_P=B(u)Q(u)(x).
\]
On montre que \(x_Q\in\Ker(Q(u))\).  
En effet :
\[
Q(u)x_Q
=
Q(u)A(u)P(u)(x)
=
A(u)P(u)Q(u)(x)
=
0.
\]
De même :
\[
P(u)x_P
=
P(u)B(u)Q(u)(x)
=
B(u)P(u)Q(u)(x)
=
0.
\]
Donc :
\[
x=x_Q+x_P\in\Ker(Q(u))+\Ker(P(u)).
\]

La somme est directe.  
Si :
\[
x\in\Ker(P(u))\cap\Ker(Q(u)),
\]
alors :
\[
x=A(u)P(u)(x)+B(u)Q(u)(x)=0.
\]
Donc l'intersection est réduite à \(\{0\}\).  
Ainsi :
\[
\Ker((PQ)(u))=\Ker(P(u))\oplus\Ker(Q(u)).
\]
\end{proof}

\subsection{Critère par polynôme annulateur scindé simple}

\begin{theorembox}
Si \(u\) admet un polynôme annulateur scindé à racines simples, alors \(u\) est diagonalisable.
\end{theorembox}

\begin{proof}
Supposons qu'il existe :
\[
P(X)=\prod_{i=1}^r(X-\lambda_i),
\]
avec \(\lambda_1,\dots,\lambda_r\) deux à deux distincts, tel que :
\[
P(u)=0.
\]
Alors :
\[
E=\Ker(P(u)).
\]
Comme les polynômes \(X-\lambda_i\) sont deux à deux premiers entre eux, le théorème de décomposition des noyaux donne :
\[
E=
\bigoplus_{i=1}^r \Ker(u-\lambda_i\Id_E).
\]
Autrement dit :
\[
E=\bigoplus_{i=1}^r E_{\lambda_i}(u).
\]
Donc \(E\) est somme directe de sous-espaces propres de \(u\).  
Ainsi \(u\) est diagonalisable.
\end{proof}

\begin{methodbox}[Diagonaliser avec un polynôme annulateur]
Pour utiliser un polynôme annulateur :
\begin{enumerate}
    \item trouver une relation polynomiale vérifiée par \(u\), par exemple \(u^2=u\), \(u^2=\Id\), \(u^3=\Id\), etc. ;
    \item écrire un polynôme \(P\) tel que \(P(u)=0\) ;
    \item factoriser \(P\) sur \(\K\) ;
    \item si \(P\) est scindé à racines simples, conclure directement que \(u\) est diagonalisable.
\end{enumerate}
\end{methodbox}

%=========================================================
\section{Projecteurs, symétries et cas classiques}

\subsection{Projecteurs}

\begin{definitionbox}[Projecteur]
Un endomorphisme \(p\in\Lcal(E)\) est un projecteur si :
\[
p^2=p.
\]
\end{definitionbox}

\begin{theorembox}
Si \(p\) est un projecteur, alors :
\[
E=\Ker(p)\oplus\Imf(p).
\]
\end{theorembox}

\begin{proof}
Pour tout \(x\in E\), on écrit :
\[
x=(x-p(x))+p(x).
\]
On a bien :
\[
p(x)\in\Imf(p).
\]
De plus :
\[
p(x-p(x))=p(x)-p^2(x)=p(x)-p(x)=0.
\]
Donc :
\[
x-p(x)\in\Ker(p).
\]
Ainsi :
\[
E=\Ker(p)+\Imf(p).
\]

Montrons que la somme est directe.  
Soit :
\[
y\in\Ker(p)\cap\Imf(p).
\]
Comme \(y\in\Imf(p)\), il existe \(x\in E\) tel que :
\[
y=p(x).
\]
Alors :
\[
p(y)=p^2(x)=p(x)=y.
\]
Mais \(y\in\Ker(p)\), donc :
\[
p(y)=0.
\]
Ainsi :
\[
y=0.
\]
La somme est donc directe.
\end{proof}

\begin{theorembox}
Tout projecteur est diagonalisable.  
Ses seules valeurs propres possibles sont \(0\) et \(1\).
\end{theorembox}

\begin{proof}
Puisque :
\[
p^2=p,
\]
on a :
\[
p^2-p=0.
\]
Donc :
\[
X(X-1)
\]
est un polynôme annulateur de \(p\).  
Ce polynôme est scindé à racines simples.  
Donc \(p\) est diagonalisable.

Si \(\lambda\) est une valeur propre de \(p\), alors elle est racine de tout polynôme annulateur, donc :
\[
\lambda(\lambda-1)=0.
\]
Ainsi :
\[
\lambda\in\{0,1\}.
\]
\end{proof}

\subsection{Symétries}

\begin{definitionbox}[Symétrie]
Un endomorphisme \(s\in\Lcal(E)\) est une symétrie si :
\[
s^2=\Id_E.
\]
\end{definitionbox}

\begin{theorembox}
Si \(s^2=\Id_E\) et si \(\operatorname{car}(\K)\neq2\), alors \(s\) est diagonalisable.  
Ses valeurs propres possibles sont :
\[
-1\quad\text{et}\quad 1.
\]
\end{theorembox}

\begin{proof}
On a :
\[
s^2-\Id_E=0.
\]
Donc :
\[
(X^2-1)(s)=0.
\]
Or :
\[
X^2-1=(X-1)(X+1).
\]
Si \(\operatorname{car}(\K)\neq2\), les racines \(1\) et \(-1\) sont distinctes.  
Le polynôme est donc scindé à racines simples.  
Ainsi \(s\) est diagonalisable.
\end{proof}

\begin{retenirbox}
\[
p^2=p \Longrightarrow p\text{ diagonalisable}.
\]
\[
s^2=\Id_E \Longrightarrow s\text{ diagonalisable si }\operatorname{car}(\K)\neq2.
\]
\end{retenirbox}

%=========================================================
\section{Trigonalisabilité}

\subsection{Définition}

\begin{definitionbox}[Endomorphisme trigonalisable]
Un endomorphisme \(u\in\Lcal(E)\) est trigonalisable s'il existe une base de \(E\) dans laquelle la matrice de \(u\) est triangulaire supérieure.
\end{definitionbox}

\begin{definitionbox}[Matrice trigonalisable]
Une matrice \(A\in\M_n(\K)\) est trigonalisable s'il existe \(P\in GL_n(\K)\) et une matrice triangulaire supérieure \(T\) telles que :
\[
A=PTP^{-1}.
\]
\end{definitionbox}

\subsection{Critère de trigonalisabilité}

\begin{theorembox}
Un endomorphisme \(u\in\Lcal(E)\) est trigonalisable sur \(\K\) si et seulement si son polynôme caractéristique est scindé sur \(\K\).
\end{theorembox}

\begin{proof}
Supposons \(u\) trigonalisable.  
Dans une certaine base, sa matrice est triangulaire supérieure :
\[
T=
\begin{pmatrix}
\lambda_1&*&\cdots&*\\
0&\lambda_2&\cdots&*\\
\vdots&\ddots&\ddots&\vdots\\
0&\cdots&0&\lambda_n
\end{pmatrix}.
\]
Alors :
\[
\chi_u(X)=\chi_T(X)=\prod_{k=1}^n(X-\lambda_k).
\]
Donc \(\chi_u\) est scindé.

Réciproquement, supposons \(\chi_u\) scindé.  
On raisonne par récurrence sur \(n=\dim(E)\).

Si \(n=1\), le résultat est évident.

Supposons \(n\geq2\).  
Comme \(\chi_u\) est scindé, il possède une racine \(\lambda\in\K\).  
Donc \(u\) possède un vecteur propre non nul \(e_1\).  
La droite :
\[
F=\Vect(e_1)
\]
est stable par \(u\), car :
\[
u(e_1)=\lambda e_1\in F.
\]
On complète \(e_1\) en une base de \(E\).  
Dans cette base, la matrice de \(u\) est de la forme :
\[
\begin{pmatrix}
\lambda&*\\
0&B
\end{pmatrix}.
\]
Le polynôme caractéristique de \(B\) est encore scindé.  
Par hypothèse de récurrence, \(B\) est trigonalisable.  
On obtient alors une base de \(E\) dans laquelle la matrice de \(u\) est triangulaire supérieure.
\end{proof}

\begin{corollaire}
Sur \(\C\), toute matrice carrée est trigonalisable.
\end{corollaire}

\begin{proof}
Sur \(\C\), tout polynôme non constant est scindé.  
En particulier, tout polynôme caractéristique est scindé.  
Donc toute matrice carrée complexe est trigonalisable.
\end{proof}

\begin{erreurbox}
Sur \(\R\), toute matrice n'est pas trigonalisable.  
Par exemple :
\[
A=
\begin{pmatrix}
0&-1\\
1&0
\end{pmatrix}
\]
a pour polynôme caractéristique :
\[
\chi_A(X)=X^2+1,
\]
qui n'est pas scindé sur \(\R\).
\end{erreurbox}

%=========================================================
\section{Sous-espaces stables}

\subsection{Définition}

\begin{definitionbox}[Sous-espace stable]
Soit \(u\in\Lcal(E)\).  
Un sous-espace vectoriel \(F\subset E\) est stable par \(u\) si :
\[
u(F)\subset F.
\]
Cela signifie :
\[
\forall x\in F,\quad u(x)\in F.
\]
\end{definitionbox}

\begin{retenirbox}
Les sous-espaces propres sont toujours stables par \(u\).  
En effet, si \(x\in E_\lambda(u)\), alors :
\[
u(x)=\lambda x\in E_\lambda(u).
\]
\end{retenirbox}

\subsection{Lien avec les matrices par blocs}

\begin{theorembox}
Soit \(F\) un sous-espace stable par \(u\), de dimension \(r\).  
Si l'on choisit une base de \(F\) que l'on complète en une base de \(E\), alors la matrice de \(u\) dans cette base est de la forme :
\[
\begin{pmatrix}
A&B\\
0&C
\end{pmatrix}.
\]
\end{theorembox}

\begin{proof}
Soit \((e_1,\dots,e_r)\) une base de \(F\), complétée en une base :
\[
(e_1,\dots,e_r,e_{r+1},\dots,e_n)
\]
de \(E\).

Pour \(j\leq r\), on a \(e_j\in F\), et comme \(F\) est stable :
\[
u(e_j)\in F.
\]
Donc les coordonnées de \(u(e_j)\) selon \(e_{r+1},\dots,e_n\) sont nulles.  
Cela signifie que le bloc inférieur gauche est nul.  
La matrice a donc la forme annoncée.
\end{proof}

%=========================================================
\section{Applications de la diagonalisation}

\subsection{Calcul des puissances}

\begin{theorembox}
Si \(A=PDP^{-1}\), avec \(D\) diagonale, alors pour tout \(n\in\N\),
\[
A^n=PD^nP^{-1}.
\]
\end{theorembox}

\begin{proof}
On raisonne par récurrence.  
Pour \(n=0\), on a :
\[
A^0=I_n=P I_n P^{-1}=PD^0P^{-1}.
\]
Pour \(n=1\), c'est l'hypothèse.

Supposons :
\[
A^n=PD^nP^{-1}.
\]
Alors :
\[
A^{n+1}=A^nA
=
(PD^nP^{-1})(PDP^{-1})
=
PD^nDP^{-1}
=
PD^{n+1}P^{-1}.
\]
\end{proof}

\subsection{Suites vectorielles}

Si une suite vectorielle est définie par :
\[
X_{n+1}=AX_n,
\]
alors :
\[
X_n=A^nX_0.
\]
Si \(A=PDP^{-1}\), alors :
\[
X_n=PD^nP^{-1}X_0.
\]

\begin{methodbox}[Résoudre une récurrence vectorielle]
Pour résoudre :
\[
X_{n+1}=AX_n,
\]
on peut :
\begin{enumerate}
    \item diagonaliser \(A\), si possible ;
    \item écrire \(A=PDP^{-1}\) ;
    \item obtenir :
    \[
    X_n=PD^nP^{-1}X_0 ;
    \]
    \item calculer explicitement \(D^n\).
\end{enumerate}
\end{methodbox}

%=========================================================
\section{Exercices progressifs}

\subsection*{Exercices du thème 1 : valeurs propres}

\begin{exercicebox}[1]
Soit :
\[
A=
\begin{pmatrix}
2&0\\
0&3
\end{pmatrix}.
\]
Déterminer les valeurs propres et les sous-espaces propres de \(A\).
\end{exercicebox}

\begin{exercicebox}[2]
Soit :
\[
A=
\begin{pmatrix}
4&1\\
0&4
\end{pmatrix}.
\]
Déterminer les valeurs propres et les sous-espaces propres de \(A\).  
La matrice est-elle diagonalisable ?
\end{exercicebox}

\subsection*{Exercices du thème 2 : polynôme caractéristique}

\begin{exercicebox}[3]
Calculer le polynôme caractéristique de :
\[
A=
\begin{pmatrix}
2&1&0\\
0&2&0\\
0&0&3
\end{pmatrix}.
\]
Déterminer ses valeurs propres.
\end{exercicebox}

\begin{exercicebox}[4]
Soit :
\[
A=
\begin{pmatrix}
0&-1\\
1&0
\end{pmatrix}.
\]
Étudier les valeurs propres de \(A\) sur \(\R\), puis sur \(\C\).
\end{exercicebox}

\subsection*{Exercices du thème 3 : diagonalisation}

\begin{exercicebox}[5]
Étudier la diagonalisabilité de :
\[
A=
\begin{pmatrix}
1&1&0\\
0&1&0\\
0&0&2
\end{pmatrix}.
\]
\end{exercicebox}

\begin{exercicebox}[6]
Diagonaliser, si possible :
\[
A=
\begin{pmatrix}
1&1\\
1&1
\end{pmatrix}.
\]
\end{exercicebox}

\subsection*{Exercices du thème 4 : polynômes annulateurs}

\begin{exercicebox}[7]
Soit \(p\in\Lcal(E)\) tel que :
\[
p^2=p.
\]
Montrer que \(p\) est diagonalisable et préciser ses valeurs propres possibles.
\end{exercicebox}

\begin{exercicebox}[8]
Soit \(s\in\Lcal(E)\) tel que :
\[
s^2=\Id_E.
\]
Montrer que \(s\) est diagonalisable si \(\operatorname{car}(\K)\neq2\).
\end{exercicebox}

\subsection*{Exercices du thème 5 : trigonalisabilité}

\begin{exercicebox}[9]
Montrer que toute matrice de \(\M_2(\C)\) est trigonalisable.
\end{exercicebox}

\begin{exercicebox}[10]
Soit :
\[
A=
\begin{pmatrix}
0&-1\\
1&0
\end{pmatrix}.
\]
Montrer que \(A\) n'est pas trigonalisable sur \(\R\), mais qu'elle l'est sur \(\C\).
\end{exercicebox}

%=========================================================
\newpage
\section{Corrigés détaillés des exercices progressifs}

\begin{correctionbox}[1]
On a :
\[
\chi_A(X)=\det
\begin{pmatrix}
X-2&0\\
0&X-3
\end{pmatrix}
=(X-2)(X-3).
\]
Les valeurs propres sont donc :
\[
2\quad\text{et}\quad3.
\]

Pour \(\lambda=2\),
\[
A-2I=
\begin{pmatrix}
0&0\\
0&1
\end{pmatrix}.
\]
Le système donne \(y=0\). Donc :
\[
E_2=\Vect
\begin{pmatrix}
1\\0
\end{pmatrix}.
\]

Pour \(\lambda=3\),
\[
A-3I=
\begin{pmatrix}
-1&0\\
0&0
\end{pmatrix}.
\]
Le système donne \(x=0\). Donc :
\[
E_3=\Vect
\begin{pmatrix}
0\\1
\end{pmatrix}.
\]
\end{correctionbox}

\begin{correctionbox}[2]
La matrice est triangulaire supérieure, donc :
\[
\chi_A(X)=(X-4)^2.
\]
L'unique valeur propre est \(4\).

On calcule :
\[
A-4I=
\begin{pmatrix}
0&1\\
0&0
\end{pmatrix}.
\]
Le système donne :
\[
y=0.
\]
Donc :
\[
E_4=\Vect
\begin{pmatrix}
1\\0
\end{pmatrix}.
\]
Ainsi :
\[
\dim(E_4)=1<2.
\]
La matrice n'est donc pas diagonalisable.
\end{correctionbox}

\begin{correctionbox}[3]
La matrice est triangulaire supérieure. Donc :
\[
\chi_A(X)=(X-2)^2(X-3).
\]
Les valeurs propres sont :
\[
2\quad\text{et}\quad3.
\]
\end{correctionbox}

\begin{correctionbox}[4]
On calcule :
\[
\chi_A(X)=\det
\begin{pmatrix}
X&1\\
-1&X
\end{pmatrix}
=X^2+1.
\]
Sur \(\R\), ce polynôme n'a pas de racine.  
Donc \(A\) n'a pas de valeur propre réelle.

Sur \(\C\),
\[
X^2+1=(X-i)(X+i).
\]
Les valeurs propres complexes sont donc :
\[
i\quad\text{et}\quad -i.
\]
\end{correctionbox}

\begin{correctionbox}[5]
La matrice est triangulaire supérieure :
\[
A=
\begin{pmatrix}
1&1&0\\
0&1&0\\
0&0&2
\end{pmatrix}.
\]
Donc :
\[
\chi_A(X)=(X-1)^2(X-2).
\]

Pour \(\lambda=1\),
\[
A-I=
\begin{pmatrix}
0&1&0\\
0&0&0\\
0&0&1
\end{pmatrix}.
\]
Le système donne :
\[
y=0,\qquad z=0.
\]
Donc :
\[
E_1=\Vect
\begin{pmatrix}
1\\0\\0
\end{pmatrix}.
\]
Ainsi :
\[
\dim(E_1)=1<2=m_1.
\]
La matrice n'est pas diagonalisable.

Pour information, pour \(\lambda=2\),
\[
A-2I=
\begin{pmatrix}
-1&1&0\\
0&-1&0\\
0&0&0
\end{pmatrix},
\]
ce qui donne :
\[
E_2=\Vect
\begin{pmatrix}
0\\0\\1
\end{pmatrix}.
\]
La somme des dimensions vaut \(1+1=2<3\).
\end{correctionbox}

\begin{correctionbox}[6]
On calcule :
\[
\chi_A(X)=\det
\begin{pmatrix}
X-1&-1\\
-1&X-1
\end{pmatrix}
=(X-1)^2-1=X^2-2X.
\]
Donc :
\[
\chi_A(X)=X(X-2).
\]
Les valeurs propres sont :
\[
0\quad\text{et}\quad2.
\]
Elles sont distinctes, donc \(A\) est diagonalisable.

Pour \(\lambda=0\),
\[
A=
\begin{pmatrix}
1&1\\
1&1
\end{pmatrix}.
\]
Le système :
\[
x+y=0
\]
donne :
\[
E_0=\Vect
\begin{pmatrix}
1\\-1
\end{pmatrix}.
\]

Pour \(\lambda=2\),
\[
A-2I=
\begin{pmatrix}
-1&1\\
1&-1
\end{pmatrix}.
\]
Le système donne :
\[
x=y.
\]
Donc :
\[
E_2=\Vect
\begin{pmatrix}
1\\1
\end{pmatrix}.
\]

On peut prendre :
\[
P=
\begin{pmatrix}
1&1\\
-1&1
\end{pmatrix},
\qquad
D=
\begin{pmatrix}
0&0\\
0&2
\end{pmatrix}.
\]
Alors :
\[
A=PDP^{-1}.
\]
\end{correctionbox}

\begin{correctionbox}[7]
Comme :
\[
p^2=p,
\]
on obtient :
\[
p^2-p=0.
\]
Donc :
\[
X^2-X=X(X-1)
\]
est un polynôme annulateur de \(p\).  
Il est scindé à racines simples.  
Donc \(p\) est diagonalisable.

Si \(\lambda\) est une valeur propre de \(p\), alors :
\[
\lambda(\lambda-1)=0.
\]
Donc :
\[
\lambda\in\{0,1\}.
\]
\end{correctionbox}

\begin{correctionbox}[8]
Comme :
\[
s^2=\Id_E,
\]
on a :
\[
s^2-\Id_E=0.
\]
Donc :
\[
X^2-1=(X-1)(X+1)
\]
est un polynôme annulateur de \(s\).

Si \(\operatorname{car}(\K)\neq2\), les racines \(1\) et \(-1\) sont distinctes.  
Le polynôme est donc scindé à racines simples.  
Donc \(s\) est diagonalisable.
\end{correctionbox}

\begin{correctionbox}[9]
Soit \(A\in\M_2(\C)\).  
Son polynôme caractéristique est un polynôme de degré \(2\) à coefficients complexes.  
Sur \(\C\), tout polynôme de degré \(2\) est scindé.  
Donc \(\chi_A\) est scindé.  
Par le critère de trigonalisabilité, \(A\) est trigonalisable.
\end{correctionbox}

\begin{correctionbox}[10]
On a déjà calculé :
\[
\chi_A(X)=X^2+1.
\]
Sur \(\R\), ce polynôme n'est pas scindé.  
Donc \(A\) n'est pas trigonalisable sur \(\R\).

Sur \(\C\),
\[
X^2+1=(X-i)(X+i).
\]
Le polynôme caractéristique est scindé.  
Donc \(A\) est trigonalisable sur \(\C\).  
Comme les deux valeurs propres sont distinctes, elle est même diagonalisable sur \(\C\).
\end{correctionbox}

%=========================================================
\newpage
\section{Grand devoir bilan final}

\begin{center}
\Large\textbf{Devoir bilan -- Réduction d'endomorphismes}
\end{center}

\subsection*{Exercice 1 -- Questions de cours et compréhension}

Soit \(E\) un espace vectoriel de dimension finie et \(u\in\Lcal(E)\).

\begin{enumerate}
    \item Définir une valeur propre et un vecteur propre.
    \item Définir le sous-espace propre \(E_\lambda(u)\).
    \item Montrer que \(E_\lambda(u)\) est un sous-espace vectoriel de \(E\).
    \item Montrer que des vecteurs propres associés à des valeurs propres deux à deux distinctes sont libres.
    \item Énoncer le critère de diagonalisation par les dimensions des sous-espaces propres.
\end{enumerate}

\subsection*{Exercice 2 -- Diagonalisation complète}

Soit :
\[
A=
\begin{pmatrix}
2&0&0\\
1&3&0\\
0&0&4
\end{pmatrix}.
\]

\begin{enumerate}
    \item Calculer \(\chi_A(X)\).
    \item Déterminer les valeurs propres de \(A\).
    \item Calculer les sous-espaces propres.
    \item Montrer que \(A\) est diagonalisable.
    \item Donner une matrice \(P\) et une matrice diagonale \(D\) telles que :
    \[
    A=PDP^{-1}.
    \]
\end{enumerate}

\subsection*{Exercice 3 -- Cas non diagonalisable}

Soit :
\[
B=
\begin{pmatrix}
1&1&0\\
0&1&1\\
0&0&1
\end{pmatrix}.
\]

\begin{enumerate}
    \item Calculer \(\chi_B(X)\).
    \item Déterminer \(E_1(B)\).
    \item La matrice \(B\) est-elle diagonalisable ?
    \item Est-elle trigonalisable ?
\end{enumerate}

\subsection*{Exercice 4 -- Polynôme annulateur}

Soit \(u\in\Lcal(E)\) tel que :
\[
u^3=u.
\]

\begin{enumerate}
    \item Donner un polynôme annulateur de \(u\).
    \item Factoriser ce polynôme.
    \item En déduire que \(u\) est diagonalisable si \(\operatorname{car}(\K)\neq2\).
    \item Quelles sont les valeurs propres possibles de \(u\) ?
\end{enumerate}

\subsection*{Exercice 5 -- Application aux suites}

Soit :
\[
A=
\begin{pmatrix}
1&1\\
1&1
\end{pmatrix}.
\]
On définit une suite de vecteurs \((X_n)\) par :
\[
X_{n+1}=AX_n,
\qquad
X_0=
\begin{pmatrix}
a\\b
\end{pmatrix}.
\]

\begin{enumerate}
    \item Diagonaliser \(A\).
    \item Calculer \(A^n\).
    \item En déduire une expression explicite de \(X_n\).
\end{enumerate}

\newpage

%=========================================================
\section{Correction détaillée du devoir bilan}

\subsection*{Correction de l'exercice 1}

\textbf{1.}  
Un scalaire \(\lambda\in\K\) est une valeur propre de \(u\) s'il existe un vecteur non nul \(x\in E\) tel que :
\[
u(x)=\lambda x.
\]
Le vecteur \(x\) est appelé vecteur propre associé à \(\lambda\).

\textbf{2.}  
Le sous-espace propre associé à \(\lambda\) est :
\[
E_\lambda(u)=\Ker(u-\lambda\Id_E).
\]

\textbf{3.}  
Comme \(u-\lambda\Id_E\) est linéaire, son noyau est un sous-espace vectoriel.  
Donc \(E_\lambda(u)\) est un sous-espace vectoriel.

\textbf{4.}  
Soient \(x_1,\dots,x_p\) des vecteurs propres associés à des valeurs propres distinctes \(\lambda_1,\dots,\lambda_p\).  
On raisonne par récurrence comme dans le cours.  
L'idée est de partir d'une relation :
\[
\alpha_1x_1+\cdots+\alpha_px_p=0,
\]
d'appliquer \(u\), puis de soustraire \(\lambda_p\) fois la première relation.  
On obtient :
\[
\alpha_1(\lambda_1-\lambda_p)x_1+\cdots+
\alpha_{p-1}(\lambda_{p-1}-\lambda_p)x_{p-1}=0.
\]
Par récurrence, tous les coefficients correspondants sont nuls.  
Puis \(\alpha_p=0\).  
Donc la famille est libre.

\textbf{5.}  
Un endomorphisme \(u\) est diagonalisable si et seulement si :
\[
\sum_{\lambda\in\Sp(u)}\dim(E_\lambda(u))=\dim(E).
\]

\subsection*{Correction de l'exercice 2}

La matrice est triangulaire inférieure.  
Donc :
\[
\chi_A(X)=(X-2)(X-3)(X-4).
\]
Les valeurs propres sont :
\[
2,\quad3,\quad4.
\]
Elles sont distinctes.  
Donc \(A\) est diagonalisable.

Calculons les sous-espaces propres.

Pour \(\lambda=2\),
\[
A-2I=
\begin{pmatrix}
0&0&0\\
1&1&0\\
0&0&2
\end{pmatrix}.
\]
Le système donne :
\[
x+y=0,\qquad z=0.
\]
Donc :
\[
E_2=\Vect
\begin{pmatrix}
1\\-1\\0
\end{pmatrix}.
\]

Pour \(\lambda=3\),
\[
A-3I=
\begin{pmatrix}
-1&0&0\\
1&0&0\\
0&0&1
\end{pmatrix}.
\]
Le système donne :
\[
x=0,\qquad z=0.
\]
Donc :
\[
E_3=\Vect
\begin{pmatrix}
0\\1\\0
\end{pmatrix}.
\]

Pour \(\lambda=4\),
\[
A-4I=
\begin{pmatrix}
-2&0&0\\
1&-1&0\\
0&0&0
\end{pmatrix}.
\]
Le système donne :
\[
x=0,\qquad y=0.
\]
Donc :
\[
E_4=\Vect
\begin{pmatrix}
0\\0\\1
\end{pmatrix}.
\]

On peut choisir :
\[
P=
\begin{pmatrix}
1&0&0\\
-1&1&0\\
0&0&1
\end{pmatrix},
\qquad
D=
\begin{pmatrix}
2&0&0\\
0&3&0\\
0&0&4
\end{pmatrix}.
\]
Alors :
\[
A=PDP^{-1}.
\]

\subsection*{Correction de l'exercice 3}

La matrice \(B\) est triangulaire supérieure. Donc :
\[
\chi_B(X)=(X-1)^3.
\]
L'unique valeur propre est \(1\), de multiplicité algébrique \(3\).

Calculons :
\[
B-I=
\begin{pmatrix}
0&1&0\\
0&0&1\\
0&0&0
\end{pmatrix}.
\]
Le système :
\[
(B-I)X=0
\]
donne :
\[
y=0,\qquad z=0.
\]
Donc :
\[
E_1(B)=\Vect
\begin{pmatrix}
1\\0\\0
\end{pmatrix}.
\]
Ainsi :
\[
\dim(E_1)=1<3.
\]
Donc \(B\) n'est pas diagonalisable.

En revanche, \(B\) est déjà triangulaire supérieure.  
Donc elle est trigonalisable.

\subsection*{Correction de l'exercice 4}

On a :
\[
u^3=u.
\]
Donc :
\[
u^3-u=0.
\]
Un polynôme annulateur est :
\[
P(X)=X^3-X.
\]
On factorise :
\[
P(X)=X(X^2-1)=X(X-1)(X+1).
\]
Si \(\operatorname{car}(\K)\neq2\), les racines :
\[
-1,\quad0,\quad1
\]
sont distinctes.  
Le polynôme annulateur est donc scindé à racines simples.  
Donc \(u\) est diagonalisable.

Les valeurs propres possibles sont les racines de \(P\), donc :
\[
\Sp(u)\subset\{-1,0,1\}.
\]

\subsection*{Correction de l'exercice 5}

On a :
\[
A=
\begin{pmatrix}
1&1\\
1&1
\end{pmatrix}.
\]
Son polynôme caractéristique est :
\[
\chi_A(X)=X(X-2).
\]
Les valeurs propres sont :
\[
0\quad\text{et}\quad2.
\]

On a :
\[
E_0=\Vect
\begin{pmatrix}
1\\-1
\end{pmatrix},
\qquad
E_2=\Vect
\begin{pmatrix}
1\\1
\end{pmatrix}.
\]

On prend :
\[
P=
\begin{pmatrix}
1&1\\
-1&1
\end{pmatrix},
\qquad
D=
\begin{pmatrix}
0&0\\
0&2
\end{pmatrix}.
\]
Alors :
\[
A=PDP^{-1}.
\]

Calculons \(P^{-1}\).  
On a :
\[
\det(P)=1\cdot1-1\cdot(-1)=2.
\]
Donc :
\[
P^{-1}
=
\frac12
\begin{pmatrix}
1&-1\\
1&1
\end{pmatrix}.
\]

Ainsi :
\[
A^n=PD^nP^{-1}.
\]
Pour \(n\geq1\),
\[
D^n=
\begin{pmatrix}
0&0\\
0&2^n
\end{pmatrix}.
\]
Donc :
\[
A^n
=
\begin{pmatrix}
1&1\\
-1&1
\end{pmatrix}
\begin{pmatrix}
0&0\\
0&2^n
\end{pmatrix}
\frac12
\begin{pmatrix}
1&-1\\
1&1
\end{pmatrix}.
\]
On obtient :
\[
A^n
=
2^{n-1}
\begin{pmatrix}
1&1\\
1&1
\end{pmatrix}.
\]
Donc :
\[
X_n=A^nX_0
=
2^{n-1}
\begin{pmatrix}
1&1\\
1&1
\end{pmatrix}
\begin{pmatrix}
a\\b
\end{pmatrix}.
\]
Ainsi, pour \(n\geq1\),
\[
X_n=
2^{n-1}
\begin{pmatrix}
a+b\\
a+b
\end{pmatrix}.
\]

%=========================================================
\newpage
\section{Fiche de synthèse finale}

\subsection*{1. Vocabulaire}

\[
u\in\Lcal(E)
\]
signifie que \(u\) est un endomorphisme de \(E\).

\[
\lambda\in\Sp(u)
\]
signifie que \(\lambda\) est une valeur propre de \(u\).

\[
u(x)=\lambda x,\qquad x\neq0
\]
signifie que \(x\) est un vecteur propre associé à \(\lambda\).

\[
E_\lambda(u)=\Ker(u-\lambda\Id_E)
\]
est le sous-espace propre associé à \(\lambda\).

\subsection*{2. Polynôme caractéristique}

Pour \(A\in\M_n(\K)\) :
\[
\chi_A(X)=\det(XI_n-A).
\]

\[
\lambda\in\Sp(A)
\Longleftrightarrow
\chi_A(\lambda)=0.
\]

Pour une matrice triangulaire, les valeurs propres sont les coefficients diagonaux.

\subsection*{3. Multiplicités}

\[
m_\lambda=\text{multiplicité algébrique de }\lambda.
\]

\[
\dim(E_\lambda)=\text{multiplicité géométrique de }\lambda.
\]

Toujours :
\[
1\leq \dim(E_\lambda)\leq m_\lambda.
\]

\subsection*{4. Familles de vecteurs propres}

Des vecteurs propres associés à des valeurs propres deux à deux distinctes sont libres.

Les sous-espaces propres associés à des valeurs propres distinctes sont en somme directe.

\subsection*{5. Diagonalisation}

\(u\) est diagonalisable si et seulement s'il existe une base de \(E\) formée de vecteurs propres.

Critère géométrique :
\[
u\text{ diagonalisable}
\Longleftrightarrow
E=\bigoplus_{\lambda\in\Sp(u)}E_\lambda(u).
\]

Critère dimensionnel :
\[
u\text{ diagonalisable}
\Longleftrightarrow
\sum_{\lambda\in\Sp(u)}\dim(E_\lambda)=\dim(E).
\]

Critère avec multiplicités :
\[
u\text{ diagonalisable}
\Longleftrightarrow
\chi_u\text{ est scindé et }
\forall\lambda,\ \dim(E_\lambda)=m_\lambda.
\]

Cas rapide :
\[
\dim(E)=n,\quad u\text{ a }n\text{ valeurs propres distinctes}
\Longrightarrow
u\text{ diagonalisable}.
\]

\subsection*{6. Matrices diagonalisables}

\[
A\text{ diagonalisable}
\Longleftrightarrow
\exists P\in GL_n(\K),\ \exists D\text{ diagonale},\quad A=PDP^{-1}.
\]

Les colonnes de \(P\) sont les vecteurs propres.

Les coefficients diagonaux de \(D\) sont les valeurs propres correspondantes.

\subsection*{7. Polynômes annulateurs}

\[
P(u)=0
\]
signifie que \(P\) est un polynôme annulateur de \(u\).

Si \(P(u)=0\) et \(\lambda\in\Sp(u)\), alors :
\[
P(\lambda)=0.
\]

Si \(u\) admet un polynôme annulateur scindé à racines simples, alors \(u\) est diagonalisable.

\subsection*{8. Projecteurs et symétries}

Projecteur :
\[
p^2=p.
\]
Alors :
\[
p\text{ diagonalisable},\qquad \Sp(p)\subset\{0,1\}.
\]

Symétrie :
\[
s^2=\Id_E.
\]
Si \(\operatorname{car}(\K)\neq2\), alors :
\[
s\text{ diagonalisable},\qquad \Sp(s)\subset\{-1,1\}.
\]

\subsection*{9. Trigonalisabilité}

\(u\) est trigonalisable si et seulement s'il existe une base dans laquelle sa matrice est triangulaire supérieure.

Critère :
\[
u\text{ trigonalisable}
\Longleftrightarrow
\chi_u\text{ est scindé sur }\K.
\]

Sur \(\C\), toute matrice carrée est trigonalisable.

\subsection*{10. Sous-espaces stables}

Un sous-espace \(F\subset E\) est stable par \(u\) si :
\[
u(F)\subset F.
\]

Si \(F\) est stable et si une base de \(F\) est complétée en une base de \(E\), alors la matrice de \(u\) est de la forme :
\[
\begin{pmatrix}
A&B\\
0&C
\end{pmatrix}.
\]

%=========================================================
\section*{Conclusion}
\addcontentsline{toc}{section}{Conclusion}

La réduction des endomorphismes est l'un des chapitres centraux de l'algèbre linéaire en classe préparatoire.  
Elle repose sur deux idées complémentaires :

\begin{itemize}
    \item une idée géométrique : chercher des droites ou sous-espaces stables, notamment les sous-espaces propres ;
    \item une idée algébrique : exploiter le polynôme caractéristique et les polynômes annulateurs.
\end{itemize}

La diagonalisation est la forme la plus efficace de réduction, mais elle n'est pas toujours possible.  
La trigonalisabilité constitue une forme plus souple, possible dès que le polynôme caractéristique est scindé.

Le point essentiel à retenir est le suivant :
\[
\text{Réduire, c'est choisir une base adaptée.}
\]

\end{document}